\begin{abstract}
最適輸送（optimal transport; OT）の理論は，フランスの数学者 Gaspard Monge（1746--1818）の言葉を借りて非形式的に
次のように述べられる：シャベルを手にした作業員が，建設現場に積まれた大きな砂の山を移動させなければならない．
作業員の目標は，その砂すべてを用いて，あらかじめ定められた形状の目標の山（たとえば巨大な砂の城）を築くことである．
当然ながら作業員は，たとえばシャベル一杯の砂を運ぶのに費やす総距離や総時間で測られる総労力を最小化したいと考える．
%
最適輸送に関心を持つ数学者は，この問題を2つの確率分布---同じ体積をもつ2つの異なる砂の山---を比較する問題として
捉える．彼らは，第1の山を第2の山へ変形・\emph{輸送}・整形する数多くのやり方のすべてを考え，砂粒1つをある場所から
別の場所へ動かすのにかかる「局所的」なコストを用いて，そのような輸送のそれぞれに「大域的」なコストを対応づける．
数学者が関心を寄せるのは，最もコストの小さい輸送のもつ性質，ならびにその効率的な計算である．
%
この最小コストは分布間の距離を定めるのみならず，確率分布の空間上に豊かな幾何構造をもたらす．この構造は，
分布が定義される台となる「基底」空間の鍵となる幾何的性質を引き継ぐという意味で標準的（canonical）である．
たとえば，基底空間がユークリッド空間であるとき，補間・重心・凸性・関数の勾配といった鍵概念が，OT の幾何を備えた
分布の空間へと自然に拡張される．

OT は多くの枠組み・多様な形で（再）発見されてきており，それゆえ豊かな歴史をもつ．Monge の先駆的研究が工学的な
問題に動機づけられていた一方で，1920 年代の Tolstoi，1940 年代の Hitchcock，Kantorovich，Koopmans は，それが
ロジスティクスや経済学において重要であることを確立した．Dantzig は 1949 年に線形計画の枠組みでこれを数値的に解き，
OT に最適化における確固たる基盤を与えた．OT はその後 1990 年代に，とくに Brenier をはじめとする解析学者たちによって
再訪され，同時に earth mover's distance の名のもとに計算機ビジョンの分野でも有名になった．
%
近年では，大きな問題次元へとスケールできる近似ソルバの登場により，OT の普及にもう一つの革命が起きている．
その結果，OT は画像科学（色やテクスチャの処理など），グラフィックス（形状操作），機械学習（回帰・分類・生成モデリング）
といった様々な問題を解き明かすために，ますます用いられるようになっている．

本稿は，数値手法に重きを置いて OT を概観し，新しいアルゴリズムの設計を導きうる OT の理論的性質を扱う．
%
とくに，OT がデータ科学において意義をもつのを助けてきた，近年の効率的アルゴリズムの潮流に焦点を当てる．
ここ数年で提案されてきた OT の数多くの一般化に重要な位置づけを与え，それらを統計的推論・カーネル法・情報理論に
由来する関連手法と結びつける．
%
本稿のすべての図は，付随するウェブサイト\footnote{\url{https://optimaltransport.github.io/}}で公開されている
コードによって再現できる．このウェブサイトでは書籍プロジェクト Computational Optimal Transport を扱っており，
スライドや計算資源も入手できる．
\end{abstract}
